\documentclass[11pt]{article}
\usepackage{amsmath,amssymb,amsthm,mathtools,mathrsfs}
\usepackage[margin=25mm]{geometry}
\usepackage[hidelinks]{hyperref}
\newcommand{\C}{\mathbb C}
\newcommand{\B}{\mathscr B}
\newcommand{\M}{\mathcal M}
\newcommand{\F}{\mathcal F_+}
\title{Finite-prime dual transfers in the original theta source}
\author{}\date{20 September 2026}
\begin{document}\maketitle

\section{Human source and exact original conventions}
Connes, Consani and Moscovici \cite[Propositions groundstate,
sonintensor and groundstatedual; Theorem isosonin]{CCM} construct two
different semilocal transfers and prove their dual-pairing identity.
We give the complete receiving calculation below on the original
positive-real Hilbert space, then on the original theta complex and
the literal tensor packet. The factors of two, prime powers, Mellin
signs and actual arithmetic mass are not discarded.

Put $\mathcal H=L^2((0,\infty),dx)$, with inner product
$\langle f,g\rangle=\int_0^\infty\overline f g\,dx$.
The cosine transform is
\[
 (\F f)(x)=2\int_0^\infty\cos(2\pi xt)f(t)\,dt.
                                                               \tag{SD1}
\]
It is a unitary involution on $\mathcal H$ by the ordinary Fourier
Plancherel theorem applied to even extensions; each even extension
has squared norm twice the positive-real norm. The original generator
is $D=-x\partial_x$. Its centered Hilbert realization is
$(D-1/2)/i$, not a replacement for $D$.
Use the left Mellin transform and its unitary critical-line restriction
\[
 \M f(s)=\int_0^\infty f(x)x^{s-1}\,dx,\qquad
 Uf(t)=\frac{\M f(1/2+it)}{\sqrt{2\pi}}.                       \tag{SD2}
\]
Indeed $x=e^y$ sends $f$ isometrically to $e^{y/2}f(e^y)$ in
$L^2(dy)$, and SD2 is its Fourier transform with the positive
exponential and factor $(2\pi)^{-1/2}$. Thus all subsequent scalar
products use $dt$, not an unrecorded multiple of it.

Let $\mathcal P$ be a fixed finite set of finite primes. For $p\in
\mathcal P$ define the unitary dilation
\[
 (V_pf)(x)=p^{1/2}f(px),\quad
 (V_p^*f)(x)=p^{-1/2}f(x/p),\quad r_p=p^{-1/2}.
                                                               \tag{SD3}
\]
Change of variables proves both the adjoint and unitarity. On their
natural domains these dilations commute with $D$, and all $V_p$ and
$V_q^*$ commute. Directly from SD1, $\F V_p=V_p^*\F$.

\section{Two bounded transfers and the preserved cross-pairing}
Set
\[
 E_{\mathcal P}=\prod_{p\in\mathcal P}(I-r_pV_p)^{-1},\qquad
 T_{\mathcal P}=\prod_{p\in\mathcal P}(I-r_pV_p^*)
                 =(E_{\mathcal P}^*)^{-1}.                     \tag{SD4}
\]
All inverse factors exist by norm-convergent geometric series since
$r_p<1$. This proves the full Hilbert-space domains and bounded
inverses, without an infinite Euler-product limit.
For one prime the formulas on the dense test space are
\[
 E_pf(x)=\sum_{j=0}^\infty f(p^jx),\qquad
 T_pf(x)=f(x)-p^{-1}f(x/p).                                    \tag{SD5}
\]
The first series is also the Hilbert norm limit because its $j$th
term has norm $p^{-j/2}\|f\|$.
Writing $s=1/2+it$, SD2 gives the exact multipliers
\[
 UE_{\mathcal P}f(t)=e_{\mathcal P}(s)Uf(t),\quad
 e_{\mathcal P}(s)=\prod_p(1-p^{-s})^{-1},
\]
\[
 UT_{\mathcal P}f(t)=b_{\mathcal P}(s)Uf(t),\quad
 b_{\mathcal P}(s)=\prod_p(1-p^{s-1}),\qquad
 \overline{b_{\mathcal P}(s)}e_{\mathcal P}(s)=1.
                                                               \tag{SD6}
\]
The plus sign in SD2 is why the two spectral arguments here are the
opposites of the paper's negative-exponential variable. No scalar
factor is inferred from the paper's other Hilbert conventions.
Equation SD4, or directly SD6 and Plancherel, proves
\[
 \boxed{\langle T_{\mathcal P}f,E_{\mathcal P}g\rangle
 =\langle f,g\rangle=\langle E_{\mathcal P}f,T_{\mathcal P}g\rangle.}
                                                               \tag{SD7}
\]
This keeps the complete complex pairing, rather than its real part.
Let
\[
 b_- =\prod_p(1-r_p),\qquad b_+=\prod_p(1+r_p).
\]
The triangle and reverse triangle inequalities in SD6 give
\[
 b_-\|f\|\le\|T_{\mathcal P}f\|\le b_+\|f\|,
 \qquad b_+^{-1}\|f\|\le\|E_{\mathcal P}f\|\le b_-^{-1}\|f\|.
                                                               \tag{SD8}
\]
These are explicit finite-set constants, not uniform assertions as
the prime set grows.

\section{The exact Sonin-space receiver}
The maps in SD4 are not interchangeable. Define the target Fourier
involution on the same Hilbert coordinate space by
\[
 \mathcal F_{+,\mathcal P}=E_{\mathcal P}\F E_{\mathcal P}^{-1}
   =E_{\mathcal P}(E_{\mathcal P}^*)^{-1}\F.                    \tag{SD9}
\]
The second equality follows by commuting $\F$ past each dilation
factor. The operator $E_{\mathcal P}$ is normal, since all factors
commute with their adjoints. Consequently
$E_{\mathcal P}(E_{\mathcal P}^*)^{-1}$ is unitary; hence so is
$\mathcal F_{+,\mathcal P}$. Similarity in the first expression proves its
square is one, and it is therefore selfadjoint. The exact
intertwining equations are
\[
 \mathcal F_{+,\mathcal P}E_{\mathcal P}=E_{\mathcal P}\F,
 \qquad\mathcal F_{+,\mathcal P}T_{\mathcal P}=T_{\mathcal P}\F.
                                                               \tag{SD10}
\]
For the second, $\F(E^*)^{-1}=E^{-1}\F$ and commutativity give
$E(E^*)^{-1}\F(E^*)^{-1}=(E^*)^{-1}\F$.

Let $P_\lambda$ denote the subspace supported in $[0,\lambda]$,
and put
\[
 \mathcal S_\lambda=(P_\lambda+\F P_\lambda)^\perp,
 \qquad
 \mathcal S_{\lambda,\mathcal P}
      =(P_\lambda+\mathcal F_{+,\mathcal P}P_\lambda)^\perp.              \tag{SD11}
\]
Both $f$ and its indicated Fourier transform vanish below $\lambda$
in the respective space. Formula SD5 shows
$E_{\mathcal P}P_\lambda\subset P_\lambda$ and shows that
$T_{\mathcal P}$ preserves a lower support gap $x<\lambda$.
Thus SD10 gives $T\mathcal S_\lambda\subset
\mathcal S_{\lambda,\mathcal P}$.
Conversely, for $h$ in the latter space write $h=Tf$, using the
bounded inverse of $T$. For $g\in P_\lambda$, SD7 gives
$\langle f,g\rangle=\langle h,Eg\rangle=0$.
Also $E\F g=\mathcal F_{+,\mathcal P}Eg$ lies in
$\mathcal F_{+,\mathcal P}P_\lambda$, so
$\langle f,\F g\rangle=\langle h,E\F g\rangle=0$.
This proves the complete correspondence
\[
 \boxed{T_{\mathcal P}:\mathcal S_\lambda
       \longrightarrow\mathcal S_{\lambda,\mathcal P}
       \text{ is a bounded isomorphism with bounded inverse}.}  \tag{SD12}
\]
It is the explicit Hilbert-coordinate receiver of the semilocal
Sonin stability theorem of \cite{CCM}. It is not claimed to be an
isometry: the exact metric factors are SD8.

\section{The same map on the original theta complex}
Retain the original strong source $\B$ consisting of smooth $f$ for
which every $\sup_{x>0}x^a|D^jf(x)|$ is finite, $a\in\mathbb R$,
$j\ge0$. Retain the original even Schwartz space
$V=\{\phi:\phi(0)=0,\int_\mathbb R\phi=0\}$ and map
\[
 \Theta\phi(x)=2\sum_{n\ge1}\phi(nx),\qquad Q=\B/\Theta V.
                                                               \tag{SD13}
\]
These are the original conventions, related to the cited orbit trace
by its explicitly retained factor two \cite{CCtrace,Meyer}.
Each finite dilation $f(x/p)$ preserves all the $\B$ seminorms:
\[
 \sup_x x^a|D^j[f(x/p)]|=p^a\sup_y y^a|D^jf(y)|.
\]
It also preserves even Schwartz functions and both vanishing
conditions in $V$. Define $T$ on $V$ by the same finite formula SD5.
Absolute convergence of the theta sum and dilation commutation give
\[
 T_{\mathcal P}\Theta=\Theta T_{\mathcal P},\qquad
 DT_{\mathcal P}=T_{\mathcal P}D.                             \tag{SD14}
\]
Thus the actual chain map is $(T|_V,T|_\B)$, inducing $\overline T$
on $Q$. The Mellin multiplier on the full strong source is the
entire exponential polynomial $b_{\mathcal P}(s)$ of SD6.
The inverse of $T$ on the Hilbert space is not asserted to preserve
$\B$, and $E$ is not asserted to be a strong-source chain map.
All uses of $E$ below remain on its proved Hilbert domain. This is
an explicit domain map, not a missing relation between the sources.

For a finite packet of nontrivial zeros $\rho$ with complete
multiplicities $m_\rho$, let $h(s)=\prod_\rho(s-\rho)^{m_\rho}$,
$g=2\xi$, and $\M F_h=g/h$. Since $0<\Re\rho<1$,
each $b_{\mathcal P}(\rho)\ne0$. In the divided-jet coordinates
$u_{\rho,j}$, multiplication by $b_{\mathcal P}$ is the exact
lower triangular matrix
\[
 (B_{\mathcal P,h})_{\rho;jl}=
 \begin{cases} b_{\mathcal P}^{(j-l)}(\rho)/(j-l)!,&j\ge l,\\
 0,&j<l.\end{cases}                                          \tag{SD15}
\]
For explicit coefficients, sum over every subset $A\subset\mathcal P$,
write $n_A=\prod_{p\in A}p$, and use $n_\varnothing=1$:
\[
 \beta_{\rho,j}:=\frac{b_{\mathcal P}^{(j)}(\rho)}{j!}
 =\frac1{j!}\sum_{A\subset\mathcal P}(-1)^{|A|}
          n_A^{\rho-1}(\log n_A)^j.                           \tag{SD16}
\]
Here $0^0=1$ for $j=0$ and the empty subset contributes zero for
$j>0$. The inverse block has coefficients
\[
 \gamma_{\rho,0}=\beta_{\rho,0}^{-1},\qquad
 \gamma_{\rho,j}=-\beta_{\rho,0}^{-1}
               \sum_{l=1}^j\beta_{\rho,l}\gamma_{\rho,j-l}.
                                                               \tag{SD17}
\]
Multiplication of the two finite Taylor series proves the inverse,
with no assumption of simple zeros. Therefore the restriction to
each original finite packet is invertible and commutes with its
original $M_s$ action. The original physical unit $j_h(g/h)$ stays
in the injection: its transformed value is
$B_{\mathcal P,h}j_h(g/h)$, not an omitted scalar.

\section{Tensor source, relation kernel and all cross phases}
Keep the original tensor degree $k$, the sum $S=s_1+\cdots+s_k$,
and the original cyclic algebra $C=\C[S]/(\chi_{h,k})$ with
injection $\eta:C\to(\C[s]/(h))^{\otimes k}$ carrying the full
physical unit. The actual source is
\[
 \mathcal V_{h,k}P=P(D_1+\cdots+D_k)(F_h^{\otimes k}),\quad
 J^{(k)}\mathcal V_{h,k}=\eta\pi_\chi.
                                                               \tag{SD18}
\]
By SD14 and SD15 the new source and actual arithmetic injection are
\[
 \mathcal V^T_{h,k}=T_{\mathcal P}^{\otimes k}\mathcal V_{h,k},
 \quad\eta^T=B_{\mathcal P,h}^{\otimes k}\eta,
 \qquad J^{(k)}\mathcal V^T_{h,k}=\eta^T\pi_\chi.              \tag{SD19}
\]
All factors of the tensor multiplier
$\prod_{i=1}^kb_{\mathcal P}(s_i)$ remain. It is not replaced by
$b_{\mathcal P}(S)$: these are different explicitly displayed
functions. Since $B_{\mathcal P,h}^{\otimes k}$ is invertible,
$\eta^T$ is injective and the polynomial relation kernel is
\emph{exactly the same} $(\chi_{h,k})$. This proves the original
cyclic quotient receiver without assuming that the new tensor
injection equals the old one.

Let $w_h(t)=|(g/h)(1/2+it)|^2/(2\pi)$, with its actual total mass
$\mu_h$. The three one-factor observations are
\[
 w_T(t)=|b_{\mathcal P}(1/2+it)|^2w_h(t),\quad
 w_E(t)=|e_{\mathcal P}(1/2+it)|^2w_h(t),\quad
 \overline b_{\mathcal P}e_{\mathcal P}w_h=w_h.                \tag{SD20}
\]
Using tensor Plancherel and Fubini, for all original polynomials $P,Q$,
\begin{align*}
\langle\mathcal V^T P,\mathcal V^T Q\rangle
 &=\int_\mathbb R\overline{P(k/2+iu)}Q(k/2+iu)
                      w_T^{*k}(u)\,du,\\
\langle E^{\otimes k}\mathcal V P,E^{\otimes k}\mathcal V Q\rangle
 &=\int_\mathbb R\overline{P(k/2+iu)}Q(k/2+iu)
                      w_E^{*k}(u)\,du,\\
\boxed{\langle\mathcal V^T P,E^{\otimes k}\mathcal V Q\rangle}
 &=\boxed{\int_\mathbb R\overline{P(k/2+iu)}Q(k/2+iu)
                       w_h^{*k}(u)\,du.}                    \tag{SD21}
\end{align*}
The exponential decay of the original $w_h$ and bounded finite-prime
multipliers justify all polynomial moments. In the cross identity,
the cancellation occurs separately at each original $t_i$, before
convolution. It does not discard phase or replace a complex pairing
by a norm.

In the original ordered monomial basis let $H_N,H_N^T,H_N^E$
be the three Gram matrices at admitted degree $N$. SD8 gives
\[
 b_-^{2k}H_N\preceq H_N^T\preceq b_+^{2k}H_N,\qquad
 b_+^{-2k}H_N\preceq H_N^E\preceq b_-^{-2k}H_N,\qquad
 H_N^{T,E}=H_N.                                               \tag{SD22}
\]
At $N\ge\deg\chi_{h,k}-1$, let $J_N$ be the original polynomial
quotient map to $C$ in its original polynomial basis. The physical
jet injection $\eta$ is kept separately, as in SD18. Minimum norm on the same
fibres gives
\[
 G_N^T=(J_N(H_N^T)^{-1}J_N^*)^{-1},\quad
 G_N^E=(J_N(H_N^E)^{-1}J_N^*)^{-1},
\]
\[
 b_-^{2k}G_N\preceq G_N^T\preceq b_+^{2k}G_N,\qquad
 b_+^{-2k}G_N\preceq G_N^E\preceq b_-^{-2k}G_N.                \tag{SD23}
\]
The same inequalities follow directly by minimizing SD22 over
each unchanged polynomial fibre. They retain all relation columns
and their Schur complement; they do not identify the two separately
minimizing lifts. In particular SD21 for equal polynomial lifts
does not assert equality of the cross-pairing of two different
minimum lifts.

For an explicit fixed physical target jet $u$ when $k=1$, the
map $\eta:C\to\C[s]/(h)$ is multiplication by the full unit
$\upsilon_h=j_h(g/h)$. Its inverse exists because each selected
zero was divided out with its complete order, so
$(g/h)(\rho)\ne0$. The transferred physical injection is $B_{\mathcal
P,h}\eta$. Thus the preimage class is
$\eta^{-1}B_{\mathcal P,h}^{-1}u$, and the transferred Gram in
the fixed physical coordinates is exactly
\[
 (B_{\mathcal P,h}\eta)^{-*}G_N^T
              (B_{\mathcal P,h}\eta)^{-1}
 =B_{\mathcal P,h}^{-*}
       (\eta^{-*}G_N^T\eta^{-1})B_{\mathcal P,h}^{-1}.
                                                               \tag{SD24}
\]
This follows by substituting that preimage into its quadratic form.
The middle matrix is the quotient Gram transported by the original
unit to physical jets; it is not identified with $G_N^T$ in the
polynomial frame. Both the semilocal multiplier and the full unit
therefore remain in the physical target calculation.
At general $k$, use the displayed injection $\eta^T$ and the inverse
of $B_{\mathcal P,h}^{\otimes k}$ on its image; no unjustified
inverse within a different cyclic image is taken.

Equations SD12--24 give the concrete literature-to-programme map:
semilocal Sonin stability, a strong theta-complex endomorphism,
all divided-jet inverses, an unchanged cyclic relation kernel,
and a preserved complete cross-Gram with explicit metric comparisons.
They do not identify a selfadjoint prolate spectrum with the actual
zeta zeros, and they assert no infinite-prime or growing-degree
estimate not proved here.

\begin{thebibliography}{9}
\bibitem{CCM} A.~Connes, C.~Consani and H.~Moscovici,
\emph{Zeta zeros and prolate wave operators: semilocal adelic operators},
\href{https://arxiv.org/abs/2310.18423v2}{arXiv:2310.18423v2},
4 May2024. Original author TeX read in full, including bibliography;
source labels groundstate, sonintensor, groundstatedual, isosonin,
dualht and propbb. Original archive and exact version recorded in
the accompanying human-source ledger. The cited book proofs of
the global adelic construction are not claimed as separately read.
\bibitem{CCtrace} A.~Connes and C.~Consani,
\emph{Hochschild homology, trace map and zeta-cycles},
\href{https://arxiv.org/abs/2207.10419v1}{arXiv:2207.10419v1},
Propositions h0coinv and propmapE1, Lemma propmapE.
\bibitem{Meyer} R.~Meyer,
\emph{A spectral interpretation for the zeros of the Riemann zeta function},
\href{https://arxiv.org/abs/math/0412277v3}{arXiv:math/0412277v3},
weighted-Schwartz source and closed-range theorem. Its original
completed factor is related to the present $g$ by the retained
factor $s(s-1)$.
\end{thebibliography}
\end{document}
